\section{Classification results}

This section presents the core results of the paper. All results are
negative or classificatory in nature: they identify which classes of
enterprise state predicates are decidable, which are undecidable, and
which admit only STOP-only decidability under incomplete observability.
All statements are formulated strictly within OC Core and
ETS.K\_EA.v1.1.

\subsection{Decidable predicates}

\begin{proposition}[Observationally invariant predicates]
Let $\varphi$ be a state predicate defined exclusively in terms of
observable signatures in $\Sigma_{\mathrm{obs}}$, i.e.\ there exists a
mapping
\[
\tilde{\varphi} : \Sigma_{\mathrm{obs}} \to \{\text{true},\text{false}\}
\]
such that $\varphi = \tilde{\varphi} \circ \Pi$.
Then $\varphi$ is decidable under incomplete observability.
\end{proposition}

\noindent
This proposition characterizes the maximal class of trivially decidable
predicates: those whose truth value depends only on what is directly
observable. No claim is made that such predicates are structurally rich,
informative, or diagnostically meaningful; only their logical
decidability is asserted.

\begin{theorem}[Decidability of STOP boundary predicates]
Any predicate $\varphi_{\mathrm{STOP}}$ that is equivalent to membership
of the enterprise continuum in a STOP region of
$\partial\Omega_{EA}$ is decidable under incomplete observability,
provided that STOP membership is invariant across each observational
equivalence class.
\end{theorem}

\noindent
This result reflects the special status of STOP in ETS: STOP corresponds
to a boundary condition whose observable consequences dominate all
internally compatible states.

\subsection{Undecidable predicates}

\begin{theorem}[Structural undecidability]
Let $\varphi$ be a state predicate that depends on at least one internal
OC/ETS element not uniquely determined by $\Sigma_{\mathrm{obs}}$
(e.g.\ internal axes, potentials, flows, cycles, or threshold margins).
Then $\varphi$ is undecidable under incomplete observability.
\end{theorem}

\noindent
In particular, predicates that depend on:
\begin{itemize}
\item internal configuration within $\Omega_{EA}$ away from
      $\partial\Omega_{EA}$;
\item relative ordering or margin of thresholds $\Theta_{EA}$ not
      observable through $\Pi$;
\item existence, stability, or coupling structure of internal cycles
      $C_{EA}$;
\item magnitudes or directions of flows $J_{EA}$ not uniquely encoded in
      observable signatures,
\end{itemize}
are undecidable in the sense defined in Section~\ref{sec:definitions}.

\begin{corollary}[Undecidability of non-terminal phase distinctions]
Predicates distinguishing between non-STOP phase outcomes (e.g.\ OK
versus INERTIA, or INERTIA versus COLLAPSE) are undecidable under
incomplete observability whenever the corresponding phase conditions
are not observationally invariant.
\end{corollary}

\noindent
This corollary formalizes the impossibility of reliably distinguishing
most non-terminal enterprise states solely from admissible
observations.

\subsection{STOP-only decidability}

\begin{theorem}[STOP-only classification]
Let $\varphi$ be a phase predicate taking values in
$\{\text{OK}, \text{INERTIA}, \text{COLLAPSE}, \text{STOP}\}$.
If STOP membership is observationally invariant, but the remaining phase
values are not, then $\varphi$ is STOP-only decidable.
\end{theorem}

\noindent
In this case:
\begin{itemize}
\item the occurrence of STOP can be certified unambiguously;
\item any non-STOP outcome remains undecidable under incomplete
      observability.
\end{itemize}

This establishes a strict diagnosability boundary: only terminal or
boundary states may be decidable, while all interior distinctions
collapse under observational equivalence.

\subsection{Exhaustiveness of the classification}

\begin{proposition}[Exhaustive trichotomy]
Every admissible enterprise state predicate falls into exactly one of
the following classes under incomplete observability:
\begin{enumerate}
\item decidable;
\item undecidable;
\item STOP-only decidable.
\end{enumerate}
\end{proposition}

\noindent
No fourth class exists within the OC/ETS formalism. In particular, there
is no notion of partial, approximate, probabilistic, or asymptotic
decidability admissible under the constraints of this paper.
